\subsection{Tipos de investigación}
\label{section:tiposinv}

\begin{itemize}
  \item \textbf{Investigación cuantitativa:} se recolectarán datos numéricos a partir de métricas definidas, los cuales servirán para determinar el rendimiento del clúster frente a diversas comparativas.
  \item \textbf{Investigación experimental:} el clúster funcionará como un entorno controlado que permitirá manipular variables independientes (e.g. número de computadoras) para determinar su efecto en una variable dependiente (e.g. tiempo de ejecución).
  \item \textbf{Investigación aplicada:} se creará el clúster como una herramienta asequible de construir debido a la utilización de computadoras antiguas, resolviendo el problema práctico de la viabilidad de su construcción.
\end{itemize}